\section{Constructive Characterization of \(k\)-Hyperedge-Connected Hypergraphs}
\label{sec:characterizing-k-edge-connected-hypergraphs}
%The splitting-off operation leads to a constructive characterization of $k$-edge-connected graphs \cite{Frank-book}.
In this section, we prove Theorem \ref{thm:k-EC-hypergraph-characterization} using Theorem \ref{thm:complete-splitting-off}. 
%use \Cref{thm:complete-splitting-off} to show a constructive characterization of $k$-hyperedge-connected hypergraphs (i.e., \Cref{thm:k-EC-hypergraph-characterization}). 
Throughout this section, we use the term \emph{hypergraph} to refer to a \emph{multi-hypergraph}.
%\knote{Throughout this section, hypergraph refers to multi-hypergraphs. }
%We remark that our constructive characterization can be obtained in polynomial time owing to the strongly polynomial time algorithm for h-splitting-off. We will also need the pinching operation in hypergraphs (recall \Cref{def:pinching-operation}) which is the inverse operation to complete h-splitting-off. The following is the main result of this section.
\iffalse
We first define a \emph{pinching} operation in hypergraphs analogous to the classical pinching operation encountered in the constructive characterization of $k$-edge-connected graphs. We then use this operation to state the main result of the section below.
\begin{definition}
    Let $H = (V, E)$ be a hypergraph and $p,k \in \Z_+$ be positive integers such that $p \leq k$. By \emph{$(k, p)$-pinching} hyperedges of $H$, we mean performing the following operation:
    \begin{enumerate}
        \item pick $p$ distinct hyperedges $e_1, \ldots, e_p \in E$,
        \item pick $p$ positive integers $t_1, \ldots, t_p \in \Z_+$ such that $\sum_{i \in [p]} t_i = k$,
        \item for $i \in [p]$, partition the hyperedge $e_i$ into $t_i$ non-empty parts $e_i = \uplus_{j \in [t_i]} f_i^j$,
        \item remove the hyperedges $e_1, \ldots e_p$ from the hypergraph $H$,
        \item add a new vertex $s$ and hyperedges $\{f_i^j \cup \{s\} : j \in [t_i], i \in [p]\}$ to the hypergraph $H$.
    \end{enumerate}
\end{definition}

\begin{theorem}\label{thm:k-EC-hypergraph-characterization}
    Let $k \in \Z_+$ be a positive integer. A hypergraph $H = (V, E)$ is $k$-hyperedge-connected if and only if $H$ can be obtained by starting from the single vertex hypergraph with no hyperedges and repeatedly applying one of the following two operations:
    \begin{enumerate}
        \item add a new hyperedge over a subset of (existing) vertices, and
        \item $(k, p)$-pinching hyperedges in the existing hypergraph for some $p \leq k$.
    \end{enumerate}
\end{theorem}
\fi
We define a hypergraph $H = (V, E)$ to be \emph{minimally} $k$-hyperedge-connected if $H-e$ is not $k$-hyperedge-connected for every $e\in E$. 
%the removal of any hyperedge results in a hypergraph that is not $k$-hyperedge-connected. 
%To show the characterization of \Cref{thm:k-EC-hypergraph-characterization}, 
We need the following observation about minimally $k$-hyperedge-connected hypergraphs. %We formally state and prove our structural observation in the next lemma.
\begin{lemma}\label{lem:minimally-k-EC-hypergraph-has-degree-k-vertex}
    Let $k \in \Z_+$ be a positive integer and $G = (V, E)$ be a minimally $k$-hyperedge-connected hypergraph with $|V| \geq 2$. Then, there exists a vertex $u \in V$ such that $d_G(u) = k$.
\end{lemma}
\begin{proof}
    Since $G$ is minimally $k$-hyperedge-connected, there exists a set $X \subseteq V$ such that $d_G(X) = k$. Let $X\subseteq V$ be an inclusion-wise minimal set such that $d_G(X)=k$. We note that if $|X| = 1$, then the claim holds. Suppose that $|X| \geq 2$.

    We first show that there exists a hyperedge $e\in E$ contained in the set $X$. By way of contradiction, suppose this is false. Let $u \in X$ be an arbitrary vertex. We note $X - \{u\} \not = \emptyset$ since $|X| \geq 2$. Furthermore, every hyperedge incident to $u$ is not contained in the set $X$, i.e. $\delta_G(u) \subseteq \delta_G(X)$. Thus, we have that $k \leq d_G(u) \leq d_G(X) = k$, where the first inequality is because $G$ is $k$-hyperedge-connected. Thus, all inequalities are equations, and we have that $d_G(u) = k$, contradicting minimality of the set $X$.

    Let $e\in E$ be a hyperedge contained in the set $X$. 
    Next, we show that %for the hyperedge $e$ contained in $X$, 
    there exists a set $Y \subseteq V$ such that $d_G(Y) = k$ and $e \in \delta_G(Y)$. 
    %(we note that this holds for every hyperedge in $H$ but we will not need the stronger statement for the purposes of this proof). 
    By way of contradiction, suppose this is false. Then, for every $Y \subseteq V$ such that $e \in \delta_G(Y)$, we have that $d_G(Y) > k$. Thus, the hypergraph $G - e$ is $k$-hyperedge-connected, contradicting the minimally $k$-hyperedge-connected property of $G$.

    We now complete the proof of the lemma. By minimality of the set $X$, we have that $Y-X\neq \emptyset$ and $(V-Y)-X\neq \emptyset$. 
    %\knote{$Y-X\neq \emptyset$ and $(V-Y)-X\neq \emptyset$. }
    Furthermore, since $e \subseteq X$ but $e \in \delta_G(Y)$, we have that $X-Y \not = \emptyset$ 
    %\knote{$X-Y\neq \emptyset$}
    and $X \cap Y \not = \emptyset$. 
    Thus, all four sets $Y-X, X-Y, X\cap Y, V-(X\cup Y)$ are non-empty.
    %\knote{Thus, all four sets $Y-X, X-Y, X\cap Y, V-(X\cup Y)$ are non-empty.} 
    Hence, we have the following: $2k = d_{G}(X) + d_G(Y) \geq d_G(X\cap Y) + d_G(X\cup Y) \geq 2k$, where the first inequality is by submodularity of the cut function and the second inequality is because $G$ is $k$-hyperedge-connected. 
    %\knote{For the second inequality to hold, we need $X\cup Y\neq V$. Why does this hold? Holds because $V-(X\cup Y)$ is also non-empty}
    Thus, all inequalities are equations and we have that $d_{G}(X\cap Y) = k$, contradicting minimality of the set $X$.
\end{proof}

We now restate and prove \Cref{thm:k-EC-hypergraph-characterization}. 
%We remark that the constructive characterization of \Cref{thm:k-EC-hypergraph-characterization} will rely on complete h-splitting-off and can be obtained in polynomial time owing to the strongly poly-time algorithm for complete h-splitting-off (\Cref{thm:complete-splitting-off}).

\thmConstructiveCharacterization*
\begin{proof}%[Proof of \Cref{thm:k-EC-hypergraph-characterization}.]
    %\knote{The reverse direction has some issue. Say we start from a single vertex hypergraph with no hyperedges, add a singleton hyperedge. The resulting hypergraph is not $k$-hyperedge-connected. }
    If $H$ is a $k$-hyperedge-connected hypergraph, then performing both operations mentioned in the theorem gives a $k$-hyperedge-connected hypergraph. Consequently, the reverse direction follows by induction on $|V| + |E|$. We prove the forward direction.  By way of contradiction, suppose the forward direction of the theorem is false. Let $G = (V, E_G)$ be a counter-example that minimizes $|V| + |E_G|$, i.e. $G$ is $k$-hyperedge-connected but cannot be obtained by applying operations (1) and (2). We note that $G$ is minimally $k$-hyperedge-connected since otherwise, there exists a hyperedge $e \in E_G$ such that $G - e$ is $k$-hyperedge-connected, and consequently, the hypergraph $G - e$ is also a counter-example to the forward direction of the theorem contradicting minimality of the counter-example $G$. Thus, by \Cref{lem:minimally-k-EC-hypergraph-has-degree-k-vertex}, there exists a vertex $u \in V$ such that $d_G(u) = k$. Let $H = (V - u, E_H)$ be the hypergraph obtained as a local connectivity preserving complete h-splitting-off at the vertex $u$ from $G$ as guaranteed to exist by \Cref{thm:complete-splitting-off}. Then, the hypergraph $H$ is $k$-hyperedge-connected. Moreover, $|V-u| + |E_H| < |V| + |E_G|$ because complete h-splitting-off at a vertex from a multi-hypergraph does not increase the total number of hyperedges. 
    %\knote{Why does this hold? We have that $|V-u|<|V|$. Moreover, $|E_G|\le |E_H|$ since complete h-splitting-off at a vertex from a multi-hypergraph does not increase the total number of hyperedges.}. 
    By minimality of the counter-example $G$, the hypergraph $H$ can be obtained by starting from the single vertex hypergraph with no hyperedges and repeatedly applying operations (1) and (2). 
    %We note that $H$ can be obtained from $G$ by operation (2) since it is the reverse operation to complete h-splitting-off 
    We note that $G$ can be obtained from $H$ by applying operation (2). This is because $H$ was obtained by complete h-splitting-off at vertex $u$ from $G$ where $d_G(u)=k$ and operation (2) is the inverse of a complete h-splitting-off operation at a vertex with degree exactly $k$. 
    %\knote{The proof is not informative regarding why we split-off at a degree $k$ vertex and not at an arbitrary vertex.}\shubhang{this is because a $(k, \cdot)$-pinching operation is the inverse of a complete h-splitting-off operation at a vertex with degree exactly $k$}. 
    Thus, $G$ is not a counter-example to the forward direction of the theorem, a contradiction.
\end{proof}

We remark that our proof of Theorem \ref{thm:k-EC-hypergraph-characterization} is constructive since it is equivalent to a proof by induction on $|V|+|E|$. 
Using the polynomial-time algorithm in Theorem \ref{thm:complete-splitting-off}, we get the following conclusion: given a $k$-hyperedge-connected hypergraph $H$, our proof gives a polynomial-time algorithm to construct a sequence of hypergraphs $H_0, H_1, H_2, \ldots, H_t$, where $H_0$ is the single vertex hypergraph with no hyperedges, $H_t=H$ and for each $i\in [t]$, the hypergraph $H_i$ is obtained from $H_{i-1}$ by either adding a new hyperedge over a subset of vertices in $H_{i-1}$ or by $(k,p)$-pinching hyperedges in $H_{i-1}$ for some positive integer $p\le k$. 